\section{Introduction}
\label{sec:intro}

Colorectal cancer (CRC) represents one of the leading causes of
cancer-related mortality worldwide, with early endoscopic detection and
resection of precancerous polyps serving as the clinical gold standard
for prevention~\cite{siegel2023colorectal, leufkens2012factors}. Among
varied morphology types, \textbf{sessile and flat polyps} (Paris
Classification Types IIa and IIb) present the highest diagnostic hazard
during colonoscopy~\cite{paris2003endoscopic}. Due to their low height
profile, irregular borders, and textural indistinguishability from
surrounding healthy mucosa, these subtle lesions are frequently missed
or inaccurately delineated by automated segmentation
algorithms~\cite{jha2020kvasir, bernal2017comparative}.

To address these subtle boundary ambiguities, deep learning architectures
have increasingly adopted \textbf{recurrent feedback mechanisms}~\cite{tomar2021fanet,
liang2015recurrent, alom2018recurrent}. A prominent exemplar is the Feature
Attention Network (FANet)~\cite{tomar2021fanet}, which introduces cross-epoch
recurrent feedback. In recurrent vision networks, intermediate or past
prediction masks ($m$) generated in preceding passes are re-injected as
auxiliary attention priors to early encoder feature representations. The
theoretical premise is compelling: by iteratively processing intermediate
hypotheses through spatial gating modules (\texttt{MixPool}), the network
should progressively sharpen ambiguous contours and prune false positives
through multi-pass self-guidance~\cite{tomar2021fanet, tomar2021ddanet}.

% -----------------------------------------------------------------------------
\subsection{The Empirical Paradox: Over-Segmentation and Loss Neutralization}
Despite theoretical promise, deploying recurrent feedback architectures in
biomedical segmentation reveals a severe, paradoxical failure mode:
\textbf{chronic false-positive over-segmentation}. In clinical colonoscopy
cohorts, error audits reveal that over $84\%$ of recurrent network errors
are concentrated in false positives---bleeding predictions far into non-lesion
mucosa (extending up to $40.7\text{ pixels}$ beyond the true histological boundary).

Standard machine learning intuition suggests rectifying false-positive bias
via \textbf{asymmetric loss functions} (such as the Tversky
loss~\cite{salehi2017tversky} or asymmetric Focal loss~\cite{yeung2022unified}),
which assign disproportionately higher penalties to false positives ($\alpha > \beta$).
Indeed, when evaluated in a purely feedforward setting (with recurrent feedback
disconnected), asymmetric Tversky loss delivers an immediate, statistically
significant reduction in false positives ($-3.25\text{ pp}$, $p = 0.005$).

However, our extensive factorial experiments uncover a baffling empirical
paradox: \textbf{the moment recurrent feedback is engaged, the beneficial
effect of asymmetric loss is completely abolished} (exhibiting a detrimental
interaction effect of $+4.17\text{ pp}$). Rather than eliminating spurious
boundaries, the recurrent feedback network traps the model in an echo chamber
of its own past mistakes, rendering modern loss engineering entirely
ineffective.

% -----------------------------------------------------------------------------
\subsection{The Discovery: The Dual-Role Pathology (The Feedback Trap)}
Through systematic representational and gradient diagnostics, we identify
the exact mechanical pathology causing this phenomenon, which we term the
\textbf{Feedback Trap}. The core failure stems from a pervasive design flaw
in recurrent feedback learning: \textbf{forcing prediction history to act
simultaneously as forward spatial guidance and as an unconstrained backward
optimization path}.

In conventional implementations, the spatial feature gating operator combines
the learnable internal attention map $\text{fmask}$ with the downsampled
recurrent mask $m_{\text{fg}}$ via a hard thresholding operator:
\begin{equation}
  \label{eq:hard_gate_intro}
  \text{keep}_{\text{hard}} =
    \max\left(\mathbb{I}(\text{fmask} > 0.5),\; m_{\text{fg}}\right).
\end{equation}
This formulation triggers two catastrophic mathematical failures:
\begin{enumerate}
  \item \textbf{Gradient Paralysis of Internal Attention:}
        Because the indicator function $\mathbb{I}(\cdot > 0.5)$ has a
        derivative of zero almost everywhere,
        $\frac{\partial \text{keep}}{\partial \text{fmask}} \equiv 0$.
        The learnable convolutional layers dedicated to extracting internal
        spatial attention receive zero direct supervision from the segmentation
        loss objective.
  \item \textbf{Toxic Recurrent Backpropagation and Error Amplification:}
        Concurrently, non-zero gradients backpropagate freely through the
        recurrent mask branch ($m_{\text{fg}}$). When the network produces an
        early false-positive hallucination, that erroneous mask is reinjected
        in the subsequent iteration. Because the network backpropagates
        through this unvalidated recursive loop, early encoder Batch
        Normalization distributions drift catastrophically ($D_{\text{KL}} = 33.63$
        at \texttt{e1.r1.bn3}), and mask gradient norms collapse by $79.6\%$.
        The network optimizes its representations to validate its own past
        predictions rather than ground-truth anatomy, permanently overpowering
        any loss-level penalty.
\end{enumerate}

% -----------------------------------------------------------------------------
\subsection{Framing: A Mechanistic Study of Recurrent Gradient Decoupling}
\textbf{Disentangling Mechanistic Principles from SOTA Chasing:}
We explicitly frame this paper as a \textbf{mechanistic deep learning study}
rather than an empirical effort to chase leaderboards. Our explicit scientific
goal is to isolate, diagnose, and resolve the dual-role structural flaw of
recurrent feedback in biomedical vision. We formulate the foundational
\textbf{Feedback Firewall} principle:
\begin{center}
  \textit{Forward spatial guidance must be strictly decoupled from backward gradient optimization.}
\end{center}
By enforcing gradient detachment on recurrent feedback tensors ($\text{detach}(m_{\text{fg}})$),
the recurrent prediction functions purely as an uncorrupted forward spatial prior,
completely preventing recursive error amplification.

\textbf{Justifying Adversarial Stress-Testing:}
To stress-test this pathology, we benchmark our method on the Kvasir-SEG
Sessile subset. Flat, poorly demarcated sessile lesions represent the ultimate
adversarial environment where the Feedback Trap---characterized by runaway
false-positive over-segmentation---is most destructive. By curing the trap in
this ambiguous regime, we demonstrate the fundamental robustness of the
architectural principle.

% -----------------------------------------------------------------------------
\subsection{Contributions of This Work}
To break the Feedback Trap without discarding the intrinsic benefits of
recurrent self-guidance, we propose \textbf{Detached Soft-OR Gating}---a
principled formulation combining continuous Boolean relaxation with the
Feedback Firewall. Our contributions are threefold:
\begin{enumerate}
  \item \textbf{Discovery and Diagnosis of the Feedback Trap:}
        We provide the first systematic diagnosis of gradient paralysis and
        loss neutralization in recurrent medical segmentation networks, proving
        that unconstrained recurrent backpropagation creates an echo chamber
        of amplified errors documented via layer-wise BatchNorm drift,
        feature cosine similarity, and gradient tracking.
  \item \textbf{The Feedback Firewall and Detached Soft-OR Formulation:}
        We formulate the Gradient Decoupling principle and realize it via
        Detached Soft-OR gating with zero learnable parameters ($\Delta \Theta = 0$).
        We prove analytically that this guarantees continuous gradient flow
        strictly proportional to background uncertainty
        ($\frac{\partial \text{keep}}{\partial \text{fmask}} = 1 - m_{\text{fg}}$)
        while strictly severing degenerate recurrent error loops
        ($\frac{\partial \text{keep}}{\partial m_{\text{fg}}} \equiv 0$).
  \item \textbf{Empirical Disentanglement, Generalization, and Universality:}
        Through a comprehensive $3 \times 2$ factorial analysis, we empirically
        prove that our method reactivates asymmetric loss optimization. A 4-way
        ablation study strictly disentangles the benefits of continuous
        relaxation from gradient detachment. On Kvasir-SEG (Sessile) across 5
        seeds, our approach cuts False Positive Rate by \textbf{$42.8\%$}
        ($4.30\% \to 2.46\%$), elevates mean Dice by \textbf{$+4.77$\,pp}
        ($p < 0.001$), and reduces variance by $66.5\%$. Crucially, external
        zero-shot cross-center evaluation on CVC-ClinicDB ($N = 612$,
        $p = 1.61 \times 10^{-15}$) proves robust out-of-distribution transfer,
        while verification on Recurrent Residual U-Net (R2U-Net) confirms
        that our Feedback Firewall is a universal stabilization principle for
        recurrent vision architectures.
\end{enumerate}
